%%============================================================
%% Section 5: Numerical Evidence for Spectral Tail Separation
%% To be inserted into the main paper before \end{document}
%%============================================================

\section{Numerical Evidence for Spectral Tail Separation}
\label{sec:numerics}

We present numerical experiments for the PSWF product family
$f_{mn} := \psi_m^{(c)}\overline{\psi_n^{(c)}}$
in the regime $N \sim c$.
The computations use the tridiagonal spectral construction
of PSWFs (Osipov--Rokhlin--Xiao \cite{Osipov2013})
and a high-resolution Gauss--Legendre quadrature
for all reference inner products.
All code is available in \texttt{numerics/pswf\_tail\_numerics.py}.

\subsection{Setup}
\label{sec:numerics-setup}

Prolate spheroidal wave functions $\psi_n^{(c)}$ on $[-1,1]$
are computed via eigenvectors of the
symmetric tridiagonal matrix $T^{(c)}$
whose entries in the Legendre basis are given by
(see \cite{Osipov2013} Ch.~3):
\begin{align*}
  T^{(c)}_{kk} &= k(k+1)
    + c^2 \frac{2k(k+1)-1}{(2k-1)(2k+3)}, \\
  T^{(c)}_{k,k+1} &= c^2
    \frac{(k+1)(k+2)}{(2k+3)\sqrt{(2k+1)(2k+5)}}.
\end{align*}
Eigenvectors yield Legendre coefficients
$\psi_n^{(c)}(x) = \sum_k a_k^{(n)} P_k^*(x)$,
where $P_k^*$ denotes the $L^2([-1,1])$-normalized Legendre polynomial.
Orthonormality was verified numerically:
$\max_{i,j \le 12}|\langle\psi_i,\psi_j\rangle - \delta_{ij}| < 10^{-12}$
for all tested parameter values.

The spectral projection $P_N$ is defined by
\[
  P_N f := \sum_{k=0}^{N-1} \langle f, \psi_k^{(c)} \rangle_{L^2}
  \psi_k^{(c)},
\]
with inner products evaluated via Gauss--Legendre quadrature
(700 nodes).
The tail norm is computed as
\[
  \|(I-P_N)f_{mn}\|^2
  = \|f_{mn}\|^2
  - \sum_{k=0}^{N-1}
  |\langle f_{mn}, \psi_k^{(c)} \rangle|^2,
\]
which is numerically more stable than direct subtraction.

\subsection{Spectral Tail Behavior}
\label{sec:numerics-tail}

Figure~\ref{fig:tail-vs-N} displays
\[
  \|(I-P_N)f_{mn}\|_{L^2([-1,1])}
\]
as a function of the projection dimension $N$,
for representative index pairs in the bulk,
off-diagonal, and edge regimes,
at $c = 50$ and $c = 100$.

The experiments exhibit a clear separation of regimes:
for $m, n \ll N$, the spectral tail is at machine-precision level
(below $10^{-13}$),
while for $m, n \approx N$ the tail remains $O(1)$.
Intermediate (off-diagonal) pairs show intermediate decay.

\begin{remark}
The near-zero tail values for bulk pairs
are consistent with strong numerical concentration
of PSWF products within $P_N$
in the bulk regime.
We emphasize that machine-precision saturation
is expected for the bulk from spectral concentration
theory (\cite{Osipov2013} Ch.~4--5)
and does not by itself constitute a proof of exponential decay;
it is included here as an illustration
of the regime structure.
\end{remark}

\subsection{Dependence on Index Separation}
\label{sec:numerics-dist}

Figure~\ref{fig:tail-vs-dist}
shows the dependence of
$\|(I-P_N)f_{mn}\|$ on the index distance $|m-n|$,
with $m$ fixed at $N/4$ and $N = c$.
A sharp transition is observed
as $n$ approaches the effective bandlimit scale $N \sim c$:
the tail remains negligible for $n \ll N$
and grows sharply as $n \to N$.
This behavior is consistent with
the qualitative structure posited in Conjecture~\ref{conj:pswf-tail}.

\subsection{Diagonal Regime Behavior}
\label{sec:numerics-diag}

Figure~\ref{fig:diag-tail}
illustrates the diagonal tail
$\|(I-P_N)\psi_m^{(c)2}\|$
as a function of $m$, with $N = c$.
A sharp transition is observed near $m \approx N$,
separating a numerically saturated bulk regime
from a non-negligible edge regime.

\subsection{Interpretation}
\label{sec:numerics-interp}

The numerical results are consistent with a
three-regime structure:
bulk ($m, n \ll N$),
off-diagonal, and edge ($m, n \approx N$) behavior.
This is the qualitative structure
underlying Conjectures~\ref{conj:pswf-tail}
and~\ref{conj:xry-stability},
in which spectral tail control
and quadrature stability are separated into distinct regimes.

We emphasize that these experiments are heuristic
and do not constitute rigorous bounds;
they are included to illustrate
the structure of the underlying spectral phenomena
and to motivate the analytical decomposition framework.

%%============================================================
%% Figures (to be placed or referenced by main paper)
%%============================================================

% Figure 1
%\begin{figure}[ht]
%  \centering
%  \includegraphics[width=\textwidth]{numerics/plot1_tail_vs_N.png}
%  \caption{Spectral tail
%    $\|(I-P_N)f_{mn}\|$ vs.\ projection dimension $N$
%    for representative index pairs
%    (bulk, off-diagonal, edge)
%    at $c=50$ (left) and $c=100$ (right).
%    Log scale. Vertical dotted line marks $N=c$.}
%  \label{fig:tail-vs-N}
%\end{figure}

% Figure 2
%\begin{figure}[ht]
%  \centering
%  \includegraphics[width=0.7\textwidth]{numerics/plot2_tail_vs_dist.png}
%  \caption{Spectral tail
%    $\|(I-P_N)f_{mn}\|$ at $N=c$
%    vs.\ index distance $|m-n|$,
%    with $m$ fixed at $N/4$,
%    for $c=50$ and $c=100$. Log scale.}
%  \label{fig:tail-vs-dist}
%\end{figure}

% Figure 3
%\begin{figure}[ht]
%  \centering
%  \includegraphics[width=0.7\textwidth]{numerics/plot3_diagonal_tail.png}
%  \caption{Diagonal tail
%    $\|(I-P_N)\psi_m^{(c)2}\|$
%    vs.\ index $m$ at $N=c$,
%    for $c=50$ and $c=100$. Log scale.
%    The sharp transition near $m \approx N$
%    separates bulk from edge regime.}
%  \label{fig:diag-tail}
%\end{figure}
